% BHU PYQ -- Manually transcribed & corrected
% Paper: MIB-301 : Algebra (Old Course)
% Exam: B.Sc. (Hons.) Semester III Examination, 2016-17

\documentclass[11pt,a4paper]{article}
\usepackage[a4paper,top=2.5cm,bottom=2.5cm,left=2.8cm,right=2.8cm,headheight=14pt]{geometry}
\usepackage{amsmath,amssymb}
\usepackage[T1]{fontenc}
\usepackage[utf8]{inputenc}
\usepackage{lmodern}
\usepackage{microtype}
\usepackage{fancyhdr}
\usepackage{enumitem}
\usepackage{hyperref}
\usepackage{lastpage}
\usepackage{parskip}

\hypersetup{colorlinks=true,urlcolor=black,linkcolor=black,
  pdftitle={MIB-301: Algebra (Old Course) -- BHU PYQ},pdfauthor={Banaras Hindu University}}

\pagestyle{fancy}
\fancyhf{}
\renewcommand{\headrulewidth}{0.4pt}
\renewcommand{\footrulewidth}{0.4pt}
\fancyhead[L]{\small   \textit{Banaras Hindu University}}
\fancyhead[R]{\small   \textit{MIB-301: Algebra (Old Course)}}
\fancyfoot[C]{\small Page  \thepage\ of \pageref{LastPage}}


\newlist{parts}{enumerate}{2}
\setlist[parts,1]{label=(\alph*),leftmargin=*,itemsep=5pt,topsep=3pt}
\setlist[parts,2]{label=(\roman*),leftmargin=*,itemsep=3pt,topsep=2pt}

\newcommand{\pts}[1]{\hfill   \textit{[#1]}}


\begin{document}


\begin{center}
  {\large\bfseries BANARAS HINDU UNIVERSITY}\\[2pt]
  {
\normalsize B.Sc. (Hons.) Semester III Examination, 2016-17}\\[6pt]
  \rule{0.8\linewidth}{0.5pt}\\[6pt]
  {\large\bfseries Paper No. MIB-301: Algebra (Old Course)}\\[4pt]
  {
\normalsize Time: 3 Hours\hfill Maximum Marks: 70}
\end{center}

\rule{\linewidth}{0.4pt}

\smallskip



\noindent   \textit{Instructions: Answer   \textbf{five} questions including Question~1 which is compulsory. Symbols have their usual meanings.}

\bigskip




\noindent   \textbf{Question 1.}   \textit{(Compulsory --- 2 marks each.)}

\begin{parts}
  \item Give an example of each of the following matrices: symmetric, skew-symmetric, Hermitian, and skew-Hermitian.
  \item If $A$ and $B$ are Hermitian matrices, show that $AB + BA$ is Hermitian and $AB - BA$ is skew-Hermitian.
  \item Define the Echelon form of a matrix. What is the rank of a $4 \times 4$ matrix whose entries are all 1?
  \item Let $*$ be the binary operation defined on the set $\mathbb{Q}_+$ of all positive rational numbers by $a * b = \frac{ab}{5}$. Find the identity element and the inverse of an element $a \in \mathbb{Q}_+$.
  \item Show that in a group $G$, if $(ab)^2 = a^2b^2$ for all $a, b \in G$, then $G$ must be abelian.
  \item Let $G$ be a group. If $a \in G$ is of order $n$, show that $o(g^{-1}ag) = n$ for any $g \in G$.
  \item If $H$ is a subgroup of a group $G$, show that $HH = H$.
  \item Let $G = \langle a \rangle$ be a cyclic group of order 13. Find the orders of $a^5$ and $a^9$.
  \item Define cyclic permutation, and even and odd permutations, and give one example of each.
  \item Show that a group of prime order is cyclic.
  \item Let $N$ be a normal subgroup of $G$ and $K$ be a subgroup of $G$. Show that $N \cap K$ is a normal subgroup of $K$.
\end{parts}

\medskip\rule{\linewidth}{0.2pt}

\medskip


\noindent   \textbf{Question 2.}

\begin{parts}
  \item Show that every square matrix can be expressed uniquely as $P + iQ$, where $P$ and $Q$ are Hermitian matrices. \pts{6}
  \item Let $A$ be a square matrix of order $n$. Show that:
  \[ A(\operatorname{adj} A) = (\operatorname{adj} A)A = |A|I_n \]
  where $I_n$ is the identity matrix of order $n$. \pts{6}
\end{parts}

\medskip\rule{\linewidth}{0.2pt}

\medskip


\noindent   \textbf{Question 3.}

\begin{parts}
  \item By using elementary transformations, find the rank of the following matrix:
  \[ A = 
\begin{pmatrix} 2 & 3 & -1 & -1 \\ 1 & -1 & -2 & -4 \\ 3 & 1 & 3 & -2 \\ 6 & 3 & 0 & -7 \end{pmatrix} \] \pts{6}
  \item Find all solutions of the following system of linear equations:
  \[ 
\begin{aligned}
    3x + 4y - z - 6w &= 0 \\
    2x + 3y + 2z - 3w &= 0 \\
    2x + y - 14z - 9w &= 0
  \end{aligned} \] \pts{6}
\end{parts}

\medskip\rule{\linewidth}{0.2pt}

\medskip


\noindent   \textbf{Question 4.}

\begin{parts}
  \item Prove that a finite semigroup $G$ in which both the cancellation laws hold is a group. Give an example in which this result fails to hold. \pts{6}
  \item Prove that the union of two subgroups of a group $G$ is a subgroup of $G$ if and only if one is contained in the other. \pts{6}
\end{parts}

\medskip\rule{\linewidth}{0.2pt}

\medskip


\noindent   \textbf{Question 5.}

\begin{parts}
  \item Express the permutation:
  \[ \sigma = 
\begin{pmatrix} 1 & 2 & 3 & 4 & 5 & 6 & 7 \\ 4 & 3 & 2 & 6 & 7 & 1 & 5 \end{pmatrix} \]
  as the product of disjoint cycles, determine whether it is even or odd, and find its order. \pts{6}
  \item Let $f : G \rightarrow G'$ be a group homomorphism. Show that $f$ is a monomorphism (one-to-one homomorphism) if and only if $\ker f = \{e\}$, where $e$ is the identity in $G$. \pts{6}
\end{parts}

\medskip\rule{\linewidth}{0.2pt}

\medskip


\noindent   \textbf{Question 6.}

\begin{parts}
  \item Prove that the order of a cyclic group $G = \langle a \rangle$ is equal to the order of its generator $a$. \pts{6}
  \item Prove that a subgroup $H$ of a group $G$ is a normal subgroup of $G$ if and only if $g^{-1}hg \in H$ for each $g \in G, h \in H$. \pts{6}
\end{parts}

\medskip\rule{\linewidth}{0.2pt}

\medskip


\noindent   \textbf{Question 7.}

\begin{parts}
  \item Let $f : G \rightarrow G'$ be an onto group homomorphism. If $H$ and $H'$ are normal subgroups of $G$ and $G'$ respectively, show that $f(H)$ and $f^{-1}(H')$ are normal subgroups of $G'$ and $G$ respectively. \pts{6}
  \item If $H$ and $K$ are two subgroups of a group $G$ such that $H$ is normal in $G$, prove that:
  \[ (HK)/H \cong K/(H \cap K) \] \pts{6}
\end{parts}

\vfill
\rule{\linewidth}{0.4pt}

\begin{center}\small
\href{https://bhu-exam.vercel.app/}{Click here to visit our website}\\[4pt]
\href{https://me-aryan.vercel.app/}{Click here to visit our contributor}\\[4pt]
\href{https://quashan-me.vercel.app/}{Click here to visit our contributor}
\end{center}

\end{document}
